\chapter*{List of Acronyms}
\addcontentsline{toc}{chapter}{List of Acronyms}

% ---------------------------------------------------------------
% Acronyms used throughout the book. Expansions are committed;
% optional one-line descriptions are marked [OPTIONAL DESCRIPTION]
% and can be removed or filled in. Sort alphabetically on commit.
% ---------------------------------------------------------------

\begin{description}
  \item[API]           Application Programming Interface.
  \item[CI/CD]         Continuous Integration / Continuous Delivery.
  \item[CPU]           Central Processing Unit.
  \item[FP16/FP32]     16-/32-bit Floating-Point numeric precision.
  \item[GPU]           Graphics Processing Unit.
  \item[HTTP]          HyperText Transfer Protocol.
  \item[INT4/INT8]     4-/8-bit Integer numeric precision.
  \item[JSON]          JavaScript Object Notation.
  \item[KV cache]      Key-Value cache (attention-layer activations
                       reused across decoding steps).
  \item[LLM]           Large Language Model.
  \item[LLMOps]        Operational practice for productionizing LLMs.
  \item[LoRA]          Low-Rank Adaptation (parameter-efficient
                       fine-tuning).
  \item[MLOps]         Machine-Learning Operations.
  \item[MoE]           Mixture-of-Experts (sparse model architecture).
  \item[MTTR]          Mean Time To Recover.
  \item[NLP]           Natural Language Processing.
  \item[OTel]          OpenTelemetry (observability standard).
  \item[PII]           Personally Identifiable Information.
  \item[QLoRA]         Quantized LoRA (LoRA over a quantized base
                       model).
  \item[RAG]           Retrieval-Augmented Generation.
  \item[REST]          Representational State Transfer (API style).
  \item[SLO]           Service-Level Objective.
  \item[SRE]           Site Reliability Engineering.
  \item[TPU]           Tensor Processing Unit.
  \item[TTFT]          Time-To-First-Token (streaming latency).
  \item[YAML]          YAML Ain't Markup Language (config format).

  % Add more as needed; review on each major revision.
\end{description}
